% Whole-proof self-check and independent programme-audit reading completed.
% Root integration review remains a separate check, recorded in its own ledger.

\subsection*{H9-T6. The variable incidence sign from integral monodromy}

The Newton test in H9-T5 is a useful way to prove that the variable
incidence factor is zero. The following argument instead determines its
sign when it is nonzero. It uses Grothendieck's integral monodromy pairing,
not a sign assigned to one $\ell$-adic scalar.

\textbf{Hypotheses.} Let $Y$ be the generic fibre of an odd-dimensional
Lefschetz pencil of fibre dimension $n=2m-1$. Suppose its middle
cohomology has an algebraic self-adjoint variable projector, and let
$B$ be the corresponding incidence factor of $P^m(Y)$. Assume:
\begin{enumerate}
\item for every $\ell\ne\operatorname{char}k$, the variable local system
      $\mathcal V_\ell$ is nonzero and absolutely irreducible after every
      finite-index restriction of geometric monodromy;
\item at one ordinary quadratic singular fibre the vanishing cycle in
      $\mathcal V_\ell$ is nonzero for every such $\ell$;
\item $B$ has its nondegenerate symmetric incidence duality and a
      normalized rational Poincar\'e cycle, as in Saito's construction.
\end{enumerate}
For the ample-section pencils considered in H9-T4, the algebraic variable
projector is $1-j^*L_h^{-1}j_*$ on the relevant odd degree. The first two
hypotheses are the precise Picard--Lefschetz monodromy hypotheses; their
standard generic-pencil instances are in SGA 7 II, XVIII, 6.7.2.

\textbf{Claim.} Under these hypotheses the rational divisor class
\[
 \ell_m=\frac{(-1)^m}{2}\Delta_B^*d_B
\]
is a rational polarization of $B$, in every characteristic. Consequently,
under the same monodromy hypotheses, H9-T4 and H9-T6 together prove the
primitive Theorem 4 for the full incidence quotient of a generic ample
section of an abelian variety, accepting 032 Corollary 8.4 for its fixed
part.

\paragraph{H9-L1: a rational square detected at almost all primes.}
If $r\in\mathbf Q^*$ is a square in $\mathbf Q_\ell$ for all but
finitely many primes $\ell$, then $r$ is a square in $\mathbf Q$.
In particular $r>0$.

\emph{Proof.} Otherwise the quadratic field $\mathbf Q(\sqrt r)$
defines a nonprincipal quadratic Dirichlet character $\chi$ of some
conductor. At primes outside the conductor, local squareness implies
$\chi(\ell)=1$, so only finitely many Euler factors of $L(s,\chi)$
differ from those of $\zeta(s)$, for real $s>1$. The quotient
$L(s,\chi)/\zeta(s)$ therefore has a finite positive nonzero limit as
$s\downarrow1$. On the other hand, a nonprincipal periodic character
has bounded partial sums and zero sum over a period. Summation by parts
shows that $\sum_{n\ge1}\chi(n)n^{-s}$ stays bounded as $s\downarrow1$
(indeed converges at $s=1$), whereas $\zeta(s)\to+\infty$.
This is a contradiction. No Riemann hypothesis is used.\hfill$\square$

\paragraph{H9-L2: the rank-one monodromy coefficient.}
Let an abelian variety have semistable reduction over a strictly
henselian trait, and suppose the toric rank is one. Let $L_0$ be a
genuine polarization and let $E_0$ be the alternating form on the Tate
module defined by $c_1(L_0)$. Choose a primitive generator $v_0$ of the
toric cocharacter lattice, and choose the same generator of
$\mathbf Q_\ell(1)$ in the monodromy operator and in the values of $E_0$.
Then the tame monodromy logarithm has the form
\[
 N_\ell(x)=-q_0 E_0(x,v_0)v_0,\qquad q_0\in\mathbf Q_{>0},
                                                        \tag{H9.12}
\]
where the rational number $q_0$ is independent of $\ell$.

\emph{Derivation from the cited theorem.} SGA 7 I, IX, 9.2, printed
pp.~419--421, identifies
the monodromy operator with the scalar extension of the integral
pairing between the two toric character lattices. Theorem 10.4(a),
printed p.~444, asserts that this integral pairing is common to all
primes; 10.4(b) asserts positivity after composing with a genuine
polarization. To specify the lattice conventions, use the underlined
$M,M'$ of IX, 5.5 and 10.4: $M$ is the character lattice of the torus
of the dual abelian variety, whereas $M'$ is the character lattice
of the original torus. (The ordinary $M,M'$ used temporarily in
IX, 9.2 occur in the opposite order.) In rank one choose generators
$a\in M$, $b\in M'$ and the dual toric cocharacter $v_0$, and write
$u=u(a,b)$ and $\lambda_t(a)=s b$ for the polarization map.
Positivity in 10.4(b) is exactly $us>0$.
For a rational lift $y$ of the corresponding quotient generator,
the canonical identifications give $N(y)=u v_0$ and
$E_0(v_0,y)=s$, so
\[
 q_0=u/s=(us)/s^2>0.
\]
This proves rationality and positivity simultaneously; it never interprets
one $\ell$-adic square as a real square. The sign in (H9.12) is the minus
sign from alternating the arguments of $E_0$. Here $E_0$ is defined
by $c_1(L_0)$, or equivalently by the biextension
$m^*L_0\otimes p_1^*L_0^{-1}\otimes p_2^*L_0^{-1}$ of
VIII, 4.12.1, rather than by an unlabelled Cartier-duality pairing.
The anticommutative comparison in VIII, 2.3.11 must be respected.
Our convention agrees with the curve
case: a one-node Jacobian has $N(y)=v_0$ with
$E_0(v_0,y)=1$. The compatibility of this curve normalization with
the general monodromy pairing is explicitly SGA 7 I, IX, 12.5 and 12.7.
Under a ramified extension of traits of index $e$, the monodromy pairing
is multiplied by $e>0$ (IX, 10.3.5, printed p.~443).

\paragraph{Proof of H9-T6.}
If $B=0$ there is nothing to prove, so suppose $B\ne0$. Work over one
finite extension of the generic field on which $B$, the normalized
Poincar\'e cycle and a genuine polarization $L_0$ are defined. Let
\[
 f_\ell:V_\ell B\longrightarrow\mathcal V_\ell(m)
\]
be the map of realizations given by this cycle. As in H9-T5,
nondegenerate incidence duality makes $f_\ell$ injective. The
irreducibility hypothesis makes it an isomorphism, for every
$\ell\ne\operatorname{char}k$.

Every geometric endomorphism of $B$ is defined over a finite extension.
Tate-module faithfulness gives
\[
 \operatorname{End}^0(B_{\overline K})\otimes\mathbf Q_\ell
   \lhook\joinrel\longrightarrow
 \operatorname{End}_{G'}(V_\ell B)
\]
after a sufficiently small finite-index monodromy subgroup $G'$.
Absolute irreducibility makes the right side $\mathbf Q_\ell$.
It follows that $\operatorname{End}^0(B_{\overline K})=\mathbf Q$.
Consequently the nondegenerate symmetric class $\ell_m$ is
\[
 \ell_m=c L_0\quad\text{in }\operatorname{NS}(B)_{\mathbf Q},
 \qquad c\in\mathbf Q^* .                              \tag{H9.13}
\]
The task is to determine the sign of this single rational number.

Let $Q_\ell$ denote the middle cup-product pairing on
$\mathcal V_\ell(m)$, with values in $\mathbf Q_\ell(1)$.
The incidence formula gives the following exact sign:
\[
 E_{\ell_m}=(-1)^{m-1}f_\ell^*Q_\ell=cE_0.            \tag{H9.14}
\]
To verify it, write the mixed class of the universal cycle as
$z=\sum_i e_i\otimes u_i$, with $e_i\in H^1(B)$ and
$u_i\in H^n(Y)(m)$. In the product $z_{13}z_{23}$ the odd-degree
class $u_i$ crosses the degree-one parameter class $e_j$, producing
a minus sign. Thus
\[
 d_B=-\sum_{i,j}Q_\ell(u_i,u_j)e_i\boxtimes e_j.
\]
Pulling back to the diagonal and dividing by two shows that the
alternating form of $\Delta^*d_B/2$ is $-f_\ell^*Q_\ell$.
Multiplying by $(-1)^m$ proves (H9.14).

At the chosen singular fibre, take a finite trait extension that makes
the abelian variety semistable and the toric lattice split. Let its
ramification index relative to the original pencil parameter be $e>0$.
The nonzero Picard--Lefschetz operator has rank one. Hence the toric rank
of $B$ is one: by IX, 9.2 and the nondegeneracy in 10.4, rank of monodromy
equals rank of the toric character lattice.

SGA 7 II, XV, 3.3.3--3.3.6, printed pp.~191--192, supplies a vanishing generator
$\delta_\ell$, normalized up to sign by the primitive quotient of the
exceptional quadric. With the same tame-character convention as in
(H9.12), its formula, in relative dimension $2m-1$, is
\[
 N_\ell(x)=e(-1)^m Q_\ell(x,\delta_\ell)\delta_\ell.
                                                        \tag{H9.15}
\]
Here twists are suppressed only after choosing a generator of
$\mathbf Q_\ell(1)$. More intrinsically, XV defines the untwisted
vanishing generator in $H^{2m-1}(Y)(m-1)$; tensoring it with that
generator puts it in $H^{2m-1}(Y)(m)$, and dividing both cup-product
values and the tame character by that same generator gives (H9.15).
Thus no unrecorded unit is introduced into the coefficient.

Put $v_\ell=f_\ell^{-1}\delta_\ell$. Since $f_\ell$ intertwines
monodromy, $v_\ell$ spans the toric line, so
$v_\ell=a_\ell v_0$ for some $a_\ell\in\mathbf Q_\ell^*$.
Combining (H9.14) and (H9.15) gives
\[
 N_\ell(x)=-e c a_\ell^2 E_0(x,v_0)v_0.
\]
Comparison with (H9.12), using $E_0$ nondegenerate, proves
\[
 a_\ell^2=\frac{q_0}{e c}\quad\text{for every }\ell\ne
 \operatorname{char}k.                                \tag{H9.16}
\]
The right side is the same nonzero rational number for all these primes.
H9-L1 therefore says it is a rational square and in particular positive.
Since $q_0>0$ and $e>0$, one obtains $c>0$. Equation (H9.13) now proves
that $\ell_m$ is a rational polarization.

For the final assertion, the fixed primitive and variable incidence
summands are orthogonal and their dualities are block diagonal.
H9-T4 polarizes the former and the argument just completed polarizes
the latter. A product of the two rational polarizations is ample,
and ampleness descends under the isogeny giving the decomposition.
This proves the claimed extension of Theorem 4.\hfill$\square$

\paragraph{Scope.}
The result concerns the geometric generic fibre under explicit monodromy
hypotheses. It does not assert positivity on every special fibre of a
pencil, where new incidence factors can appear. Nor does it identify
the height of a merely homologically trivial generic cycle with a
N\'eron--Tate height on $P^m$: the $E^4$ domain obstruction remains.
The original H9(b) therefore does not follow from this extension.
